\begin{proof}[Proof of first part of Theorem~\ref{onlyacm}]
Let $b$ be a Busemann function associated to $\xi$ and ${X=\nabla b}$.
Then $|X|=1$ and $\diver X=h>0$.
Furthermore, $\nabla X$ is positive definite on the orthogonal complement of $X$,
that is, in the direction to the horospheres in $M$ with center $\xi$.

Let $Y$ be a smooth vector field on $M$ with compact support such that $Y\perp X$.
Then there are constants $C_1,C_2>0$ such that $\nabla X\ge C_1$ on the support of $Y$ and perpendicularly to $X$
and such that $|Y|^2\le C_2$.

Now $C^\infty_{\rm c}(M)$ is a core for $\Delta$.
Since $\diver X=h>0$ is constant, \eqref{x1} and \eqref{x2a} yield, for $u\in C^\infty_{\rm c}(M)$,
\begin{align*}
 \int_M (\Delta u-\lambda u) P_Xu
 &= \int_M \langle\nabla_{\nabla u}X,\nabla u\rangle \ge C_1 \int_M|\nabla^\bot u|^2 \\
 &\ge \frac{C_1}{C_2}\int_M|Yu|^2 = \frac{C_1}{C_2}\|T_Yu\|_2^2,
\end{align*}
where $\nabla^\bot u$ denotes the component of $\nabla u$ perpendicular to $X$.
Now \eqref{x3} applies and shows that $T_Y$ is $S$-smooth on any given bounded Borel subset $A\subseteq\R$.
Hence \eqref{Kato3} implies that $\ran E_AT_Y^*\subseteq H_{\ac}(\Delta)$,
for any bounded Borel subset $A\subseteq\R$.
Therefore, $\ran T_Y^*\subseteq H_{\ac}(\Delta)$.

Now suppose that $u\in L^2(D )$ is perpendicular to $H_{\ac}(\Delta)$.
Then $u$ is perpendicular to $\ran T_Y^*$, for all vector fields $Y$ as above.
But then $T_Yu=0$ in the sense of distributions, for all such vector fields $Y$.
This implies that $u$ is constant along horospheres, that is, levels of $u$ are unions of horospheres.
Now the preimage $b^{-1}(B)$ of any Borel subset $B\subseteq\R$ of positive measure has infinite measure.
Therefore, $u$ vanishes, by square-integrability.
\end{proof}


\begin{proof}[Proof of part of Theorem~\ref{onlyacn}]
By the definition of elementary groups of isometries of $M$, the gradient field $X$ of Busemann functions,
that is, the field of velocity vectors of unit speed geodesics emanating from $\xi$,
is invariant under $\Gamma$.
Notice that $X$ spans the normal space to the fibers of the Riemannian submersion $\pi$.
Now the arguments of the previous proof apply and show that $u\in L^2(N)$ is perpendicular to $H_{\ac}(N)$
if and only if $u$ is constant on the fibers of $\pi$.
There are now two cases:
The fibers of $\pi$ have infinite volume.
This is always the case when $\Gamma$ is of type \ref{ax}.
Then, as in the previous proof, the square-integrability of $u$ implies that $u=0$.

In the second case, $\Gamma$ is of type \ref{pa}.
Then the flow $(F_t)_{t\in\R}$ of $X$ induces diffeomorphisms between the fibers of $\pi=b$,
and the volume element of horospheres is multiplied by $\exp(ht)$ under $F_t$.
Hence, since the fibers of $b$ have finite volume, their volumes satisfy
\begin{align*}
	\exp(ht) \big|b^{-1}(s)\big| = \exp(hs) \big|b^{-1}(t)\big|.
\end{align*}
Furthermore, the space $L^2_b(N)$ of functions on $N$, which are constant fiberwise,
and its orthogonal complement $L^2_0(N)$ are invariant under $\Delta$.
That is, the spectrum of $\Delta$ is the combination of the spectra of $\Delta$ on $L^2_b(N)$
and of $\Delta$ on $L^2_0(N)$.
By what we said above, if $u\in L^2_0(N)$ is perpendicular to $H_{\ac}(N)$, then $u=0$.
Thus $L^2_0(N)\subseteq H_{\ac}(N)$, and we are left with discussing $\Delta$ on $L^2_b(N)$.

Note that $u\in L^2(N)$ is in $L^2_b(N)$ if and only if $u$ is a pull-back $b^*v$ for some function $v$ on $\R$.
We endow $\R$ with the measure $\mu=h_0\exp(ht){\rm d}t$,
where $h_0=\big|b^{-1}(0)\big|$ and get that
\begin{align*}%\label{unitary}
	b^* \colon \ L^2(\R,\mu) \to L^2_b(N)
\end{align*}
is a unitary transformation.
With respect to $b^*$, the Laplacian on $N$ corresponds to a diffusion operator on $\R$,
\begin{align*}%\label{diffusion}
	\Delta b^*v = b^*\bigl(- v'' - hv'\bigr).
\end{align*}
Let $\vf$ be the square root of $h_0\exp(ht)$.
Then
\begin{align*}%\label{diffusion2}
	\vf' = \frac{h}2\vf \qquad\text{and}\qquad \vf'' = \frac{h^2}4\vf.
\end{align*}
Therefore, under the unitary transformation
\begin{align*}
	L^2(\R,\mu) \to L^2(\R), \qquad v \mapsto \vf v,
\end{align*}
we obtain
\begin{align*}
	(\vf v)'' = \vf''v +2\vf'v' + \vf v'' = \vf\biggl(\frac{h^2}4 v + hv' + v''\biggr).
\end{align*}
We conclude that the Schr\"odinger operator $Sw=-w''+h^2w/4$ in $L^2(\R)$
corresponds to the above diffusion operator on $L^2(\R,\mu)$.
Since $h^2/4$ is a constant, the spectrum of the latter is absolutely continuous and equal to $\bigl[h^2/4,\infty\bigr)$.
By unitary equivalence, the same holds for $\Delta$ on~$L^2_b(N)$.

By Observation~\ref{chee}, $\sigma(N)\subseteq\bigl[h^2/4,\infty\bigr)$, hence the above shows equality,
$\sigma(N)=\bigl[h^2/4,\infty\bigr)$, in the case where the fibers of $\pi$ are of finite volume.
The case where the fibers are of infinite volume will be discussed in Section~\ref{secsig}.
\end{proof}
